\documentclass[conference]{IEEEtran}
\IEEEoverridecommandlockouts
\usepackage{cite}
\usepackage{amsmath,amssymb,amsfonts}
\usepackage{algorithmic}
\usepackage{graphicx}
\usepackage{textcomp}
\usepackage{xcolor}
\usepackage{booktabs}
\usepackage{multirow}
\usepackage{array}
\usepackage[hidelinks]{hyperref}

\def\BibTeX{{\rm B\kern-.05em{\sc i\kern-.025em b}\kern-.08em
    T\kern-.1667em\lower.7ex\hbox{E}\kern-.125emX}}

% ───────────────────────────────────────────────────────────────
\begin{document}

\title{Comparative Study of Transfer Learning Architectures\\
for 120-Class Handwritten Javanese Script Recognition\\
Under Severe Class Imbalance}

\author{%
\parbox{\linewidth}{\centering
  \begin{minipage}[t]{.31\linewidth}\centering
    \normalsize\textbf{1\textsuperscript{st} Fathan Fardian Sanum}\\[4pt]
    \normalsize\textit{School of Computing}\\
    \textit{Telkom University}\\
    Bandung, Indonesia\\
    \small fathansanum@student.\linebreak telkomuniversity.ac.id
  \end{minipage}%
  \hfill%
  \begin{minipage}[t]{.31\linewidth}\centering
    \normalsize\textbf{2\textsuperscript{nd} Syahdan Rizqi Ruhendy}\\[4pt]
    \normalsize\textit{School of Computing}\\
    \textit{Telkom University}\\
    Bandung, Indonesia\\
    \small syahdanruhendy@student.\linebreak telkomuniversity.ac.id
  \end{minipage}%
  \hfill%
  \begin{minipage}[t]{.31\linewidth}\centering
    \normalsize\textbf{3\textsuperscript{rd} Ardia Alif Ramadhan}\\[4pt]
    \normalsize\textit{School of Computing}\\
    \textit{Telkom University}\\
    Bandung, Indonesia\\
    \small ardiaramadhan@student.\linebreak telkomuniversity.ac.id
  \end{minipage}%
}%
}

\maketitle

% ───────────────────────────────────────────────────────────────
\begin{abstract}
Javanese script (\textit{Hanacaraka}) is a traditional writing system
facing the threat of functional extinction due to a declining number of
generations able to read and write it, making digitization through
automated character recognition an urgent preservation strategy.
Three research gaps not yet addressed simultaneously in the literature
are full 120-class recognition under natural class imbalance, with an
imbalance ratio of 33.10$\times$ (largest class 695 images versus
smallest class 21 images) between the Vowel A classes and the five
other vowel groups, structurally safe augmentation that preserves
Javanese script characteristics, and comparative evaluation across
transfer learning architectures.
This study proposes a four-step structure-preserving data preparation
pipeline that explicitly avoids geometric augmentation to maintain the
visual integrity of the script, combined with a weighted loss function
to compensate for class imbalance.
Four transfer learning architectures (EfficientNet-B0,
MobileNetV3-Large, ResNet-18, and VGG-16) are evaluated
comparatively using an identical pipeline, with metrics reported
separately for majority and minority classes.
EfficientNet-B0 achieves 99.06\% accuracy with a macro F1-score of
0.9823, surpassing prior studies using artificially balanced datasets,
while simultaneously maintaining an F1 gap between majority and minority
classes of only 0.0291, the best generalization balance among all
architectures tested.
These results demonstrate that extreme class imbalance can be overcome
with a targeted pipeline, while per-vowel-group analysis shows that the
E-taling group is consistently the hardest to recognize across all four
architectures, a pattern hidden by aggregate metrics.
\end{abstract}

\begin{IEEEkeywords}
\textit{Abugida}, augmented character recognition, class-imbalanced
dataset, EfficientNet-B0, handwritten character recognition,
Javanese script, MobileNetV3-Large, ResNet-18,
structure-preserving augmentation, \textit{transfer learning}, VGG-16,
\textit{weighted loss}
\end{IEEEkeywords}

% ───────────────────────────────────────────────────────────────
\section{Introduction}

Javanese script (\textit{Hanacaraka}) is an \textit{abugida} writing
system combining 20 base consonants with six vowel markers into
120 syllabic classes~\cite{wicaksono2024}.
The dominance of Latin script in modern education has caused a
cross-generational decline in proficiency~\cite{faizin2025,wisesty2026},
making digitization via \textit{Optical Character Recognition}
(OCR) an urgent preservation intervention~\cite{widiarti2026}.
The primary challenge is that the only public dataset covering all
120 classes, \textit{Indonesian Local Script Characters}~\cite{wisesty2026},
exhibits extreme imbalance, where Vowel A averages 500 images per class
while 100 non-A classes have only around 21 images each, yielding
an imbalance ratio of 33.10$\times$.

\textit{Transfer learning} based on CNNs has proven effective for
handwritten regional script recognition with limited
data~\cite{ifriza2026,nindya2025}, significantly outperforming
\textit{shallow learning} methods.
Combined with a weighted loss function, this approach provides a
suitable framework for simultaneously addressing the visual complexity
of Javanese script and its class imbalance.

Prior studies leave three gaps not yet addressed simultaneously.
Susanto et al.~\cite{susantoIJAI2023} achieved 99.65\% accuracy on
120 classes but on an artificially balanced dataset, while Wisesty
et al.~\cite{wisesty2026} found that geometric augmentation actually
degrades Javanese script recognition performance.
These three gaps are 120-class recognition under natural class
imbalance, structurally safe augmentation strategies, and comparative
cross-architecture evaluation.

This study closes all three gaps by proposing a four-step
structure-preserving data preparation pipeline and evaluating four
architectures (EfficientNet-B0, MobileNetV3-Large, ResNet-18, VGG-16)
comparatively.
EfficientNet-B0 achieves 99.06\% accuracy together with the smallest
F1 gap across vowel groups of only 0.0291, while per-vowel-group
analysis reveals that the E-taling group is the largest source of
error across all architectures.


% ───────────────────────────────────────────────────────────────
\section{Related Work}
\label{sec:related}

The following studies are the most closely related to this work
because they directly target the full 120 syllabic class coverage
of Javanese script, serving as the primary benchmarks to surpass.
Susanto et al.~\cite{susantoIJECE2023} developed a custom CNN with
four convolutional layers and two \textit{fully connected} layers for
120-class classification, achieving 97.29\% accuracy on a balanced
dataset of 14,400 images from 30 writers.
The same group~\cite{susantoIJAI2023} subsequently proposed a
12-layer CNN with selective augmentation, where rotation was restricted
to $\pm3^{\circ}$ to avoid altering the visual meaning of the script,
achieving 99.65\% accuracy.
Despite the high accuracy figures, all experiments were conducted on
datasets artificially balanced with an equal number of samples per
class, so model performance under naturally imbalanced data
distributions has never been explored.

Beyond studies specifically targeting 120 classes, several studies
covering the broader Nusantara script landscape provide valuable
insights into architecture selection, augmentation behavior, and model
capability limits on 20-class Javanese script.
Wisesty et al.~\cite{wisesty2026} proposed MNetNCR based on
MobileNetV3~\cite{howard2019mobilenetv3} for five Nusantara scripts
and achieved F1\,=\,0.9788 for Javanese script without any augmentation.
More importantly, applying augmentation to this model actually
degraded performance to 0.9579 because geometric transformations
destroyed the visual characteristics distinguishing between classes,
serving as the primary empirical foundation for the
augmentation strategy in this study.
Nindya et al.~\cite{nindya2025} applied VGG-16~\cite{simonyan2015vgg}
to multi-script classification of Balinese, Javanese, and Sundanese
scripts, achieving 93.19\% accuracy for the Javanese subset, and
introduced the \textit{Character Likeness Rate} (CLR) metric
identifying the ``ha''/``la'' pair (CLR 25.3\%) and ``sa''/``da''
(CLR 15.8\%) as the most confusion-prone characters due to
visual similarity.
Sulistiyo et al.~\cite{sulistiyo2025} compared various
\textit{shallow learning} methods on the Javanese script dataset and
found that even the best-performing XGBoost model struggled
significantly with visually similar characters such as ``tha''
(F1-score 13.64\%) and ``nga'' (F1-score 29.63\%), demonstrating
the inherent limitations of feature-based approaches for
fine-grained character discrimination.
Fateah et al.~\cite{ifriza2026} demonstrated that ResNet-18~%
\cite{he2016resnet} fine-tuned with a \textit{transfer learning}
strategy can achieve 91\% accuracy and 91\% F1-score from only
1,562 images of Javanese script, making it the sole ResNet-based \textit{transfer learning}
reference explicitly targeting Javanese script recognition
prior to this study.
All studies in this section, while informative, are limited to
balanced datasets and have never tested model robustness under
extreme class imbalance.

Class imbalance is a critical aspect overlooked by almost all
Javanese script studies, yet it is the natural characteristic of
the publicly available dataset.
Faizin et al.~\cite{faizin2025} is the only study to explicitly
address this issue, developing SkelBAGAN
(\textit{Skeleton-based Balancing GAN}) for a Javanese script dataset
with imbalance ratios up to 81.78$\times$.
SkelBAGAN integrates a character skeleton-based \textit{autoencoder}
and \textit{adversarial loss} to generate synthetic images preserving
script structure integrity, and demonstrably improves recognition
performance on minority classes.
Its limitation lies in the GAN architecture complexity requiring
significant computational resources, and its scope remaining limited
to 20 base classes without testing on the full 120-class syllabic
scenario.
This study fills that gap with a lighter approach, namely structure-preserving
augmentation and a weighted loss function, without requiring
additional generative network training.

To understand the positioning of this study in the broader landscape,
a review of OCR studies across Southeast Asian abugida scripts
provides context on how far research in this area has advanced.
Wicaksono et al.~\cite{wicaksono2024} conducted a systematic
literature review on six Southeast Asian abugida scripts and found
no studies in the past five years addressing word segmentation for
either Javanese or Sundanese script, confirming that isolated character
recognition remains the frontier of Javanese OCR research.
Widiarti et al.~\cite{widiarti2026} established a baseline
evaluation for printed (not handwritten) Javanese OCR, noting that
handwriting style variation adds a far more complex layer of
challenges compared to printed script.
Together, these two studies underscore that the domain of
handwritten Javanese character recognition, particularly at
120-class coverage with imbalanced data, remains very open and
serves as the direct motivation for this study.

% ───────────────────────────────────────────────────────────────
\section{Dataset and Experimental Setup}
\label{sec:dataset}

This study uses the \textit{Indonesian Local Script Characters} dataset,
curated by Aditya Firman Ihsan et al.\ at Telkom University and
available on Mendeley Data, which was also used by prior
studies~\cite{wisesty2026}.
For the Javanese script subset, the dataset consists of two parts with
a total of 13,257 images, comprising the \texttt{all\_class/} subdirectory
containing 12,191 images across 120 syllabic classes (6 vowels
$\times$ 20 consonants), and the \texttt{variations/} subdirectory
containing 1,066 handwriting style variation images for 8 consonants
(ba, ga, ha, ka, la, nya, ra, ta) in the Vowel A group.
Images predominantly use JPEG format (91\%) with RGB color channels
(98.3\%) and an original average resolution of $362\times347$ pixels,
subsequently standardized to $224\times224$ pixels as model input.

The class distribution in this dataset is severely imbalanced by
nature, as summarized in Table~\ref{tab:dist}.
The twenty Vowel A classes have 426--695 images per class
($\mu = 499.8$), while 100 non-A classes have only 21--22 images
per class, yielding an imbalance ratio of 33.10$\times$
(largest/smallest class = 695/21) with a per-class median of only 22,
far below the mean of 101.59.
This imbalance is the primary challenge that all pipeline stages are
designed to address.

\begin{table}[t]
\caption{Class Distribution of Dataset by Vowel Group}
\label{tab:dist}
\begin{center}
\renewcommand{\arraystretch}{1.2}
\begin{tabular}{lccc}
\toprule
\textbf{Vowel Group} & \textbf{Classes} &
  \textbf{Images/Class} & \textbf{Total} \\
\midrule
Vowel A        & 20 & 426--695 & 9,995 \\
Vowel E        & 20 & 21--22   &   440 \\
Vowel I        & 20 & 21--22   &   440 \\
Vowel O        & 20 & 21--22   &   440 \\
Vowel U        & 20 & 21--22   &   440 \\
Vowel E-taling & 20 & 21--22   &   440 \\
\midrule
\textbf{Total} & \textbf{120} & 21--695 & \textbf{12,191} \\
\bottomrule
\multicolumn{4}{l}{\footnotesize IR = 33.10$\times$, median = 22.}
\end{tabular}
\end{center}
\end{table}

The dataset was split per-class using stratified sampling at a
70:15:15 ratio into training (8,483 original images), validation
(1,789 images), and test (1,919 images) sets.
Per-class stratification ensures all 120 classes are represented
in every subset, including minority classes with only 21 images.
After augmentation on the training set, the total training samples
increased to 38,510 images.
All images were converted to grayscale and replicated across three
channels to maintain compatibility with ImageNet-pretrained backbones,
then normalized using ImageNet statistics
($\mu\,{=}\,[0.485;\ 0.456;\ 0.406]$,
$\sigma\,{=}\,[0.229;\ 0.224;\ 0.225]$).

% ───────────────────────────────────────────────────────────────
\section{Proposed Method}
\label{sec:method}

\subsection{Pipeline Overview}

The primary challenge in this study lies not in the model architecture,
but in the data conditions, specifically the extreme class imbalance
between Vowel A and non-A classes that causes conventional training
strategies to produce models systematically biased toward majority
classes.
To address this, the study proposes a four-step data preparation
pipeline executed sequentially before the training process begins.
Each step is designed to handle one specific aspect of the data
distribution problem, namely maintaining minority class representation in
the split, adding data diversity without new collection, filling
minority classes with structurally safe augmentation, and applying
proportional penalty weights during training.
Fig.~\ref{fig:pipeline} illustrates the overall flow of this pipeline.

\begin{figure}[t]
\centerline{\includegraphics[width=\columnwidth]{inggris.drawio.png}}
\caption{Four-step data preparation pipeline applied sequentially
  before model training.}
\label{fig:pipeline}
\end{figure}

\subsection{Step 1: Per-Class Stratified Split}

The dataset is split into training (70\%), validation (15\%), and
test (15\%) sets using per-class stratified sampling.
This ensures all 120 classes are represented in every subset,
even for minority classes with only 21 images, yielding
$\approx$15 training images, 3 validation images, and 3 test images
per minority class before augmentation.
This split yields 1,789 validation images and 1,919 test images.

\subsection{Step 2: Variation Data Integration}

A total of 1,066 images from \texttt{variations/} are added to the
training set for the 8 corresponding Vowel A consonant classes
(ba, ga, ha, ka, la, nya, ra, ta).
This step increases handwriting style diversity for those classes
without requiring new data collection.

\subsection{Step 3: Structure-Preserving Augmentation}

This is the primary methodological contribution of the study.
Augmentation is applied \textit{exclusively to the 100 non-A classes}
until each class reaches a target of $T = 300$ training images.
New images are created by repeatedly sampling from existing training
images and applying transformations until each class reaches $T$
samples.
The value $T = 300$ was selected to yield approximately a $14\times$
increase over the original $\approx$21 minority images per class,
representing a deliberate balance between expanding data diversity
and limiting near-duplicate samples generated from a small source pool.

The most critical design constraint is the prohibition of geometric
transformations in any form.
Wisesty et al.~\cite{wisesty2026} empirically demonstrated that
augmentation actually degrades Javanese script recognition because
transformations alter visual characteristics that distinguish between
classes.
This is reinforced by the finding of Nindya et al.~\cite{nindya2025}
that the ``ha'' and ``la'' pair are nearly identical horizontal
reflections, so applying \textit{horizontal flip} would directly
create cross-class ambiguous labels.

Table~\ref{tab:aug} lists the permitted and prohibited transformations.
Permitted transformations target \textit{appearance variability}
(ink density, scan noise, pen stroke thickness) without altering the
spatial layout of the character.
Each augmented image is created by applying 1--3 transformations
selected randomly from the permitted list.
Augmentation is performed \textit{offline} (saved to disk) to ensure
reproducibility.
The final training set after augmentation totals 38,510 images,
an increase of approximately $4.0\times$ from 9,549 images before
augmentation (8,483 original training images plus 1,066 variation
images).

\begin{table}[t]
\caption{Structure-Preserving Augmentation Strategy}
\label{tab:aug}
\begin{center}
\renewcommand{\arraystretch}{1.2}
\begin{tabular}{p{1.7cm}p{3.6cm}c}
\toprule
\textbf{Category} & \textbf{Transformation} & \textbf{Status} \\
\midrule
\multirow{5}{1.7cm}{Photometric}
  & Brightness jitter ($\pm$30\%)                 & \checkmark \\
  & Contrast jitter ($\pm$30\%)                   & \checkmark \\
  & Gaussian blur (radius 0.5--1.5\,px)           & \checkmark \\
  & Gaussian noise ($\sigma$ = 1--5\%)            & \checkmark \\
  & Sharpening ($1.2\times$--$2.0\times$)         & \checkmark \\
\midrule
\multirow{2}{1.7cm}{Morphological}
  & Dilation (MaxFilter $3\!\times\!3$ / $5\!\times\!5$) & \checkmark \\
  & Erosion (MinFilter $3\!\times\!3$ / $5\!\times\!5$)  & \checkmark \\
\midrule
Spatial & Random crop + resize (max 10\% margin) & \checkmark \\
\midrule
\multirow{4}{1.7cm}{Geometric (prohibited)}
  & Rotation ($>\pm3^{\circ}$)          & $\times$ \\
  & Horizontal / vertical flip          & $\times$ \\
  & Shear / affine transform            & $\times$ \\
  & Elastic distortion                  & $\times$ \\
\bottomrule
\end{tabular}
\end{center}
\end{table}

\subsection{Step 4: Weighted Loss Function}

After augmentation, class weights are computed using:
\begin{equation}
  w_c = \frac{N}{C \cdot n_c}
  \label{eq:weight}
\end{equation}
where $N$ is the total training samples, $C = 120$ is the number
of classes, and $n_c$ is the number of training samples for class $c$.
This formula is equivalent to \texttt{compute\_class\_weight('balanced')}
in scikit-learn, which assigns greater loss penalties to
misclassification of minority classes.
These weights are applied as a parameter to the
\texttt{CrossEntropyLoss} function during training.

\subsection{Model Architectures}

Once the data pipeline is established, the next stage is selecting
the model architectures to evaluate.
Four CNN architectures with ImageNet-1K pretrained weights are
selected based on two complementary criteria.
The first criterion is representation of conceptually distinct
architecture families, including compound scaling efficiency (EfficientNet),
residual connections (ResNet), \textit{depthwise separable
convolution} (MobileNetV3), and deep convolutional blocks (VGG).
The second criterion is direct relevance to prior Javanese script
research, enabling fair comparison against architectures already
reported in the literature.
The smallest variant of each family was deliberately chosen so that
all models could be trained on a single 8\,GB GPU at the same input
resolution, allowing performance differences to be attributed to
architectural character rather than computational budget.
EfficientNet-B0 was selected as a representative of compound scaling
that balances depth, width, and resolution at a small parameter count.
ResNet-18 was chosen because it is the backbone reported by Fateah et
al.~\cite{ifriza2026} for Javanese script, whereas MobileNetV3-Large
follows the choice of Wisesty et al.~\cite{wisesty2026} and VGG-16
follows Nindya et al.~\cite{nindya2025}, so all three enable direct
comparison with already published results.
Table~\ref{tab:arch} summarizes the four architectures and their
literature relevance.
Each architecture's classification head is replaced with a linear
layer mapping to 120 classes, while all backbone weights are
initialized from ImageNet-1K pretraining.

\begin{table}[t]
\caption{CNN Architectures Compared}
\label{tab:arch}
\begin{center}
\renewcommand{\arraystretch}{1.2}
\begin{tabular}{lcl}
\toprule
\textbf{Architecture} & \textbf{Parameters} & \textbf{Literature Relevance} \\
\midrule
EfficientNet-B0~\cite{tan2019efficientnet}    & $\sim$5.3M  & -- \\
MobileNetV3-Large~\cite{howard2019mobilenetv3}& $\sim$5.4M  & Wisesty et al.~\cite{wisesty2026} \\
ResNet-18~\cite{he2016resnet}                 & $\sim$11.7M & Fateah et al.~\cite{ifriza2026} \\
VGG-16~\cite{simonyan2015vgg}                 & $\sim$138M  & Nindya et al.~\cite{nindya2025} \\
\bottomrule
\end{tabular}
\end{center}
\end{table}

\subsection{Training Strategy}

All models follow a uniform two-phase fine-tuning strategy to ensure
fair cross-architecture comparison.
The split into two phases is designed to prevent the large gradients
from the still-random classification head from corrupting the already
informative pretrained weights.
In Phase 1 over 10 epochs, all backbone parameters are frozen and
only the classification head is trained with a higher learning rate.
This phase aims to build a stable head initialization before gradients
propagate to the more sensitive pretrained feature layers, and for the
AdamW-based models a three-epoch \textit{warmup} is also used so that
the learning-rate transition proceeds smoothly.
In Phase 2 running up to 90 epochs, all layers are unfrozen for
full fine-tuning with a lower learning rate.
The optimizer is matched to each architecture, namely AdamW for
EfficientNet-B0 and MobileNetV3-Large which are more stable under
adaptive learning rates, and momentum SGD for ResNet-18 and VGG-16
which yields better generalization on those two networks.
During Phase 2, a ReduceLROnPlateau scheduler with factor 0.5
adaptively reduces the learning rate when validation loss stagnates,
while \textit{early stopping} with \textit{patience} 20 halts training
if no improvement occurs, and \textit{label smoothing} of 0.1 is
applied to all models to curb overconfidence toward majority classes.
All models are trained using the weighted \texttt{CrossEntropyLoss}
per~\eqref{eq:weight}.
Full hyperparameter configuration per architecture is shown in
Table~\ref{tab:hyperparam}.

\begin{table}[t]
\caption{Hyperparameter Configuration per Architecture}
\label{tab:hyperparam}
\begin{center}
\renewcommand{\arraystretch}{1.2}
\begin{tabular}{lccccc}
\toprule
\textbf{Architecture} & \textbf{Opt.} & \textbf{LR Phase 1} & \textbf{LR Phase 2}
  & \textbf{WD} & \textbf{ES} \\
\midrule
EfficientNet-B0    & AdamW & $10^{-3}$          & $10^{-4}$          & $10^{-4}$ & 20 \\
MobileNetV3-Large  & AdamW & $10^{-3}$          & $8{\times}10^{-5}$ & $10^{-5}$ & 20 \\
ResNet-18          & SGD   & $2{\times}10^{-2}$ & $10^{-3}$          & $5{\times}10^{-4}$ & 20 \\
VGG-16             & SGD   & $10^{-3}$          & $10^{-4}$          & $5{\times}10^{-4}$ & 20 \\
\bottomrule
\multicolumn{6}{l}{\footnotesize Opt. = optimizer. WD = weight decay. ES = early stopping patience.} \\
\multicolumn{6}{l}{\footnotesize Label smoothing 0.1 and batch 32 on all models, except} \\
\multicolumn{6}{l}{\footnotesize VGG-16 uses batch 16, dropout 0.5, and grad clipping 2.0.}
\end{tabular}
\end{center}
\end{table}

All experiments were run on an NVIDIA GeForce RTX 3070 Ti GPU
(8\,GB VRAM) with an AMD Ryzen 5 3600 CPU, using Python 3.10,
PyTorch 2.9.1 (CUDA 13.0), torchvision 0.24.1, and scikit-learn.
All random seeds were fixed at 42 for reproducibility, with a batch
size of 32 for most models and 16 for VGG-16, and 3 DataLoader workers.

Evaluation is performed on the test set (1,919 samples) using
Top-1 accuracy and macro F1-score as overall metrics.
Macro F1-score is chosen as the primary metric because it assigns
equal weight to each class regardless of majority class dominance.
As the primary analytical contribution, macro F1-score is also
computed separately for the Vowel A group (20 majority classes,
1,518 test samples) and the non-A vowel group (100
minority classes, 401 test samples), with the absolute difference
between the two used to measure model generalization balance.
As an indicator of statistical robustness, the 95\% confidence interval
for the macro F1-score is also computed via \textit{bootstrap} on the
test-set predictions.

% ───────────────────────────────────────────────────────────────
\section{Results and Discussion}
\label{sec:results}

\subsection{Overall Performance Comparison}

Table~\ref{tab:results} summarizes the performance of all architectures
on the test set.
EfficientNet-B0 ranks first on every metric with 99.06\% accuracy and
a macro F1-score of 0.9823 (95\% confidence interval 0.9640--0.9929),
followed by MobileNetV3-Large with 98.75\% accuracy and a macro
F1-score of 0.9663.
ResNet-18 ranks third with 97.71\% accuracy and a macro F1-score of
0.9321, while VGG-16 closes the ranking with 94.89\% accuracy and a
macro F1-score of 0.8326.
This pattern is notable because the top performers EfficientNet-B0 and
MobileNetV3-Large are the lightest architectures with about 5 million
parameters, whereas VGG-16 with its 138 million parameters ranks last,
so parameter count alone is not deterministic of performance under
these imbalanced conditions.

\begin{table*}[t]
\caption{Performance Comparison and Gap Analysis of Four Architectures on the 120-Class Test Set}
\label{tab:results}
\begin{center}
\renewcommand{\arraystretch}{1.3}
\small
\begin{tabular}{lcccccc}
\toprule
\textbf{Architecture} & \textbf{Accuracy (\%)} & \textbf{Macro F1}
  & \textbf{Vowel A F1} & \textbf{Non-A F1}
  & \textbf{Gap} & \textbf{Epochs} \\
\midrule
EfficientNet-B0   & \textbf{99.06} & \textbf{0.9823} & \textbf{0.9939} & \textbf{0.9648} & \textbf{0.0291} & 49 \\
MobileNetV3-Large & 98.75          & 0.9663          & 0.9935          & 0.9367          & 0.0568          & 50 \\
ResNet-18         & 97.71          & 0.9321          & 0.9010          & 0.8593          & 0.0416          & 51 \\
VGG-16            & 94.89          & 0.8326          & 0.9386          & 0.7232          & 0.2154          & 66 \\
\bottomrule
\multicolumn{7}{l}{\footnotesize Gap = $|$Vowel A F1 $-$ Non-A F1$|$.}
\end{tabular}
\end{center}
\end{table*}

\subsection{Vowel A vs.\ Non-A Vowel Gap Analysis}

The Gap column in Table~\ref{tab:results} reveals a pattern hidden
by aggregate metrics, which constitutes one of the primary analytical
contributions of this study.

EfficientNet-B0 not only leads on the overall metrics but also
maintains the smallest gap of 0.0291, with a Vowel A F1 of 0.9939
against a non-A Vowel F1 of 0.9648.
This shows that the structure-preserving augmentation pipeline and the
weighted loss function suppress majority-class bias almost entirely, so
the model still recognizes minority classes that each have only about
300 augmented training images.
MobileNetV3-Large follows with a comparable pattern, namely a gap of
0.0568, consistent with its strength in the study of Wisesty et
al.~\cite{wisesty2026} and confirming that a lightweight architecture
based on \textit{depthwise separable convolution} remains competitive
under the far more demanding 120-class condition.

ResNet-18 displays a qualitatively different pattern.
Its gap is indeed small (0.0416), but this stems from a relatively low
Vowel A F1 (0.9010) rather than a high non-A Vowel F1, so its balance
comes from majority-class performance dropping as well rather than from
strengthening of the minority classes.
This observation underscores the importance of reading the gap together
with both of its underlying values, because a small gap does not always
indicate good generalization.

VGG-16 exhibits the largest gap of 0.2154, with its non-A Vowel F1
falling to 0.7232 even though its Vowel A F1 remains high.
This reflects a deep convolutional architecture that tends to
over-optimize majority classes and is the most sensitive to minority
data scarcity.

The per-vowel-group analysis in Table~\ref{tab:per_vowel} deepens this
finding by decomposing the macro F1-score for the six vowel groups
separately.
EfficientNet-B0 leads or ties on almost every group, even reaching a
perfect F1 of 1.0000 on the Vowel I group.
The most consistent pattern is that the E-taling group is the hardest
across all four architectures, with the lowest F1 throughout, ranging
from 0.8732 on EfficientNet-B0 down to only 0.6159 on VGG-16.
The visual similarity of the E-taling marker to several other vowel
markers, together with its scarce sample count, is the likely cause,
and this becomes a clear target for improvement in future work.

\begin{table}[t]
\caption{Macro F1-score per Vowel Group on the Test Set}
\label{tab:per_vowel}
\begin{center}
\renewcommand{\arraystretch}{1.3}
\small
\begin{tabular}{lcccc}
\toprule
\textbf{Vowel} & \textbf{Efficient} & \textbf{Mobile}
  & \textbf{ResNet} & \textbf{VGG} \\
\textbf{Group} & \textbf{Net-B0} & \textbf{NetV3-L}
  & \textbf{-18} & \textbf{-16} \\
\midrule
Vowel A        & \textbf{0.9939} & 0.9935 & 0.9010 & 0.9386 \\
Vowel E        & \textbf{0.9889} & 0.9762 & 0.8806 & 0.6383 \\
Vowel I        & \textbf{1.0000} & 0.9746 & 0.9214 & 0.8976 \\
Vowel O        & \textbf{0.9746} & 0.9378 & 0.9126 & 0.5673 \\
Vowel U        & \textbf{0.9873} & 0.9216 & 0.8710 & 0.8969 \\
Vowel E-taling & \textbf{0.8732} & 0.8732 & 0.7111 & 0.6159 \\
\bottomrule
\multicolumn{5}{l}{\footnotesize Highest value per row is in bold.} \\
\multicolumn{5}{l}{\footnotesize E-taling is consistently the hardest group.}
\end{tabular}
\end{center}
\end{table}

\subsection{Training Curve Analysis}

Fig.~\ref{fig:curves_eff} and~\ref{fig:curves_mob} display the training
curves of the two best architectures.
The Phase 1 to Phase 2 transition (epoch 11) is clearly visible as a
sharp loss drop, confirming the effectiveness of the two-phase
fine-tuning strategy.
EfficientNet-B0 reaches its best weights at epoch 49 with a validation
accuracy of 98.99\%, while MobileNetV3-Large reaches its best weights
at epoch 50 with a validation accuracy of 98.66\%.
Both curves show a tight gap between training and validation accuracy,
indicating no meaningful overfitting despite the training set being
enriched through augmentation.
Training curves for ResNet-18 and VGG-16 are omitted as both converged
without anomalies, and their epoch counts and full performance metrics
are captured in Table~\ref{tab:results}.

\begin{figure}[t]
\centerline{\includegraphics[width=\columnwidth]%
  {../New_Model/results_efficientnet_b0/training_curves.png}}
\caption{EfficientNet-B0 training curves. The Phase 1-to-Phase 2
  transition is visible at epoch 11 and the best weights are reached
  at epoch 49.}
\label{fig:curves_eff}
\end{figure}

\begin{figure}[t]
\centerline{\includegraphics[width=\columnwidth]%
  {../New_Model/results_mobilenetv3_large/training_curves.png}}
\caption{MobileNetV3-Large training curves as the second-best
  architecture, with best weights at epoch 50.}
\label{fig:curves_mob}
\end{figure}

\subsection{Prediction Result Visualization}

To complement the quantitative analysis, Fig.~\ref{fig:samples}
through~\ref{fig:confmat} visualize the prediction behavior of
EfficientNet-B0 as the best architecture on the test set.
Fig.~\ref{fig:samples} presents example predictions on several test
images, showing that the model correctly recognizes a range of
handwriting styles on most samples, including minority classes.
Fig.~\ref{fig:pervowel} summarizes the F1-score distribution per vowel
group and reaffirms E-taling as the lowest-performing group.
Fig.~\ref{fig:confmat} shows a confusion matrix dominated by its
diagonal, with the small share of errors concentrated on visually
similar character pairs, consistent with the character-similarity
findings in prior literature.

\begin{figure}[t]
\centerline{\includegraphics[width=\columnwidth]%
  {../New_Model/results_efficientnet_b0/sample_predictions.png}}
\caption{Example predictions of EfficientNet-B0 on the test set.
  The true label and predicted label are shown for each sample.}
\label{fig:samples}
\end{figure}

\begin{figure}[t]
\centerline{\includegraphics[width=\columnwidth]%
  {../New_Model/results_efficientnet_b0/f1_per_vowel.png}}
\caption{F1-score per vowel group for EfficientNet-B0.
  The E-taling group obtains the lowest score.}
\label{fig:pervowel}
\end{figure}

\begin{figure}[t]
\centerline{\includegraphics[width=\columnwidth]%
  {../New_Model/results_efficientnet_b0/confusion_matrix.png}}
\caption{Confusion matrix of EfficientNet-B0 over 120 classes.
  The dominant diagonal indicates a low error rate.}
\label{fig:confmat}
\end{figure}

\subsection{Comparison with Prior Work}

Table~\ref{tab:comparison} contextualizes this study within the
existing Javanese script recognition literature.
This study is the first to evaluate all 120 syllabic classes on a
naturally imbalanced dataset, and the first to report four-architecture
comparisons with per-vowel-group analysis.

\begin{table}[t]
\caption{Comparison with Javanese Script Recognition Studies}
\label{tab:comparison}
\begin{center}
\renewcommand{\arraystretch}{1.2}
\begin{tabular}{p{2.6cm}ccc}
\toprule
\textbf{Study} & \textbf{Classes} & \textbf{Balanced?}
  & \textbf{Best Metric} \\
\midrule
Susanto et al.~\cite{susantoIJECE2023} & 120 & Yes & Accuracy 97.29\% \\
Susanto et al.~\cite{susantoIJAI2023}  & 120 & Yes & Accuracy 99.65\% \\
Fateah et al.~\cite{ifriza2026}        & 20  & Yes & Accuracy 91\%   \\
Wisesty et al.~\cite{wisesty2026}      & 20  & Yes & F1\,=\,0.9788  \\
Nindya et al.~\cite{nindya2025}        & 20  & Yes & Accuracy 93.19\% \\
Faizin et al.~\cite{faizin2025}        & 20  & No  & -- \\
\midrule
\textbf{This work} & \textbf{120} & \textbf{No} & \textbf{Accuracy 99.06\%} \\
\bottomrule
\end{tabular}
\end{center}
\end{table}

EfficientNet-B0 achieves 99.06\% accuracy, surpassing prior balanced
studies (97.29\% by Susanto et al.~\cite{susantoIJECE2023}) and on par
with the highest reported 99.65\% on an artificially balanced
dataset~\cite{susantoIJAI2023}, despite being trained on a dataset with
33.10$\times$ imbalance.
This confirms that the proposed pipeline successfully bridges naturally
imbalanced conditions to match the performance of ideal balanced
conditions.
Furthermore, no prior study has reported per-vowel-group analysis,
making the finding that the best generalization balance is achieved by
the lightest architecture, and that the E-taling group is consistently
the hardest, an entirely novel contribution to this literature.

One limitation must be acknowledged, as each non-A Vowel class is
represented by only 4 images in the test set, so a single misprediction
already shifts per-class F1 by 25\%.
Therefore, all conclusions are based on per-vowel-group macro F1
(averaged over 20 classes), and larger minority data collection
remains a priority for follow-up research.

% ───────────────────────────────────────────────────────────────
\section{Conclusion}
\label{sec:conclusion}

This study demonstrates that 120-class handwritten Javanese script
recognition under severe natural class imbalance can be achieved at
competitive performance levels through a combination of a
structure-preserving data preparation pipeline and transfer learning.
EfficientNet-B0 achieves 99.06\% accuracy, surpassing prior studies
using artificially balanced datasets~\cite{susantoIJECE2023}, proving
that imbalanced conditions are not an insurmountable obstacle given
an appropriate data strategy.

The analytically most significant finding is that the best
generalization balance is achieved by the lightest architecture, where
EfficientNet-B0 maintains an F1 gap between majority and minority
classes of only 0.0291 while remaining the leader on every aggregate
metric.
The per-vowel-group analysis, never reported in the Javanese script
literature before, shows that the E-taling group is consistently the
hardest to recognize across all four architectures.
This finding has direct practical implications, namely when a Javanese
OCR system is deployed on real documents spanning the full vowel
spectrum uniformly, EfficientNet-B0 is the most appropriate choice
because it combines the highest accuracy with the smallest bias toward
majority classes, while dedicated handling of the E-taling group should
become a priority for improvement.

This study is limited to isolated character recognition and a small
test set for minority classes.
Going forward, collecting more minority data will strengthen the
validation of these findings, while exploring GAN-based
augmentation~\cite{faizin2025} can complement the proposed photometric
approach.
Extension to word-level recognition incorporating \textit{pasangan}
forms represents the most application-relevant next
step~\cite{wicaksono2024}.

% ───────────────────────────────────────────────────────────────
\begin{thebibliography}{00}

\bibitem{wicaksono2024}
B. A. Wicaksono, B. S. Hantono, and T. B. Adji,
``Word segmentation task for Southeast Asian abugida scripts:
A systematic literature review,''
in \textit{Proc. 2024 2nd Int. Conf. Technology Innovation and
Its Applications (ICTIIA)}, 2024.

\bibitem{faizin2025}
M. A. Faizin, N. Suciati, and C. Fatichah,
``Handling imbalance in Javanese manuscript character dataset
using skeleton-based balancing generative adversarial networks,''
\textit{Jurnal RESTI (Rekayasa Sistem dan Teknologi Informasi)},
vol.~9, no.~4, pp.~820--831, 2025.

\bibitem{wisesty2026}
U. N. Wisesty et al.,
``MNetNCR: MobileNet model for efficient traditional Nusantara
script character recognition,''
\textit{International Journal of Artificial Intelligence (IJ-AI)},
vol.~15, no.~2, pp.~1513--1528, 2026.

\bibitem{widiarti2026}
A. R. Widiarti, S. E. P. Adji, F. T. Adji, G. K. B. Indraputra,
and Y. A. Trisnanto,
``A baseline evaluation of OCR segmentation and classification
methods for printed Javanese script,''
\textit{Engineering, Technology \& Applied Science Research (ETASR)},
vol.~16, no.~1, pp.~31699--31705, 2026.

\bibitem{ifriza2026}
N. Fateah, Subhan, and Y. N. Ifriza,
``Optimizing Javanese script recognition using fine-tuned ResNet-18
and transfer learning,''
\textit{International Journal of Artificial Intelligence (IJ-AI)},
vol.~15, no.~1, pp.~443--453, 2026.

\bibitem{nindya2025}
T. P. Nindya, M. D. Sulistiyo, and U. N. Wisesty,
``Recognition of Balinese, Javanese, and Sundanese scripts
using CNN with VGG-16 architecture,''
in \textit{Proc. 2025 Int. Conf. Data Science and Its Applications
(ICoDSA)}, 2025, pp.~1221--1226.

\bibitem{sulistiyo2025}
M. D. Sulistiyo et al.,
``Investigating shallow learning methods for optical character
recognition of Indonesia's Nusantara scripts,''
\textit{Jurnal RESTI (Rekayasa Sistem dan Teknologi Informasi)},
vol.~9, no.~6, pp.~1538--1549, 2025.

\bibitem{susantoIJAI2023}
A. Susanto et al.,
``Handwritten Javanese script recognition method based 12-layers
deep convolutional neural network and data augmentation,''
\textit{International Journal of Artificial Intelligence (IJ-AI)},
vol.~12, no.~3, pp.~1448--1458, 2023.

\bibitem{susantoIJECE2023}
A. Susanto, I. U. W. Mulyono, C. A. Sari, E. H. Rachmawanto,
and D. R. I. M. Setiadi,
``Improved Javanese script recognition using custom model of
convolution neural network,''
\textit{International Journal of Electrical and Computer Engineering
(IJECE)}, vol.~13, no.~6, pp.~6629--6636, 2023.

\bibitem{howard2019mobilenetv3}
A. Howard et al.,
``Searching for MobileNetV3,''
in \textit{Proc. IEEE/CVF Int. Conf. Computer Vision (ICCV)},
2019, pp.~1314--1324.

\bibitem{simonyan2015vgg}
K. Simonyan and A. Zisserman,
``Very deep convolutional networks for large-scale image
recognition,''
in \textit{Proc. Int. Conf. Learning Representations (ICLR)},
2015.

\bibitem{he2016resnet}
K. He, X. Zhang, S. Ren, and J. Sun,
``Deep residual learning for image recognition,''
in \textit{Proc. IEEE Conf. Computer Vision and Pattern Recognition
(CVPR)}, 2016, pp.~770--778.

\bibitem{tan2019efficientnet}
M. Tan and Q. V. Le,
``EfficientNet: Rethinking model scaling for convolutional neural
networks,''
in \textit{Proc. Int. Conf. Machine Learning (ICML)},
2019, pp.~6105--6114.

\end{thebibliography}

\end{document}
